%============================================================
% VARIABEL IDENTITAS DOKUMEN — SKRIPSI UNIVERSITAS AMIKOM PURWOKERTO
% Edit bagian ini sesuai dengan data skripsi Anda.
% File ini di-input paling awal agar semua file lain dapat
% mereferensikan variabel-variabel berikut.
%============================================================

% ── Identitas Mahasiswa ─────────────────────────────────────
% Nama: Capitalize Each Word (huruf awal tiap kata kapital).
% NIM:  ditulis dengan format titik sesuai kampus, contoh 20.21.0001.
\newcommand{\NamaMahasiswa}{Nama Mahasiswa}
\newcommand{\NIM}{xx.xx.xxxx}

% ── Judul Skripsi ────────────────────────────────────────────
% Ditulis HURUF BESAR dan tidak disingkat, maksimal 20 kata.
% Jika lebih dari 1 baris, atur pergantian baris dengan `\\`
% agar membentuk piramida terbalik (huruf V).
% Uppercase diterapkan otomatis di tempat yang membutuhkannya.
\newcommand{\JudulSkripsi}{Judul Skripsi Ditulis Huruf Besar\\
Dan Tidak Disingkat, Maksimal Dua Puluh Kata}

% Studi kasus (opsional). Kosongkan jika tidak ada.
% Contoh: (Studi Kasus: PT. Sinergi Dimensi Informatika)
\newcommand{\JudulStudiKasus}{(Studi Kasus: Jika Ada)}

% ── Dosen Pembimbing (gelar terbaru, NIDN tidak disingkat) ──
\newcommand{\NamaPembimbingSatu}{Nama Dosen Pembimbing 1, S.Kom., M.Kom.}
\newcommand{\NIDNPembimbingSatu}{00xx0xxxxxxxx}
\newcommand{\NamaPembimbingDua}{Nama Dosen Pembimbing 2, S.Kom., M.Kom.}
\newcommand{\NIDNPembimbingDua}{00xx0xxxxxxxx}

% ── Dosen Penguji (gelar terbaru, NIDN tidak disingkat) ─────
\newcommand{\NamaPengujiSatu}{Nama Penguji 1, S.Kom., M.Kom.}
\newcommand{\NIDNPengujiSatu}{00xx0xxxxxxxx}
\newcommand{\NamaPengujiDua}{Nama Penguji 2, S.Kom., M.Kom.}
\newcommand{\NIDNPengujiDua}{00xx0xxxxxxxx}

% Penguji 3 / Pendamping — diaktifkan lewat config/toggles.tex
\newcommand{\NamaPengujiTiga}{Nama Penguji 3 / Pendamping, S.Kom., M.Kom.}
\newcommand{\NIDNPengujiTiga}{00xx0xxxxxxxx}

% ── Ketua Program Studi (gelar + NIK) ────────────────────────
\newcommand{\NamaKaprodi}{Nama Ketua Program Studi, S.Kom., M.Kom.}
\newcommand{\NIKKaprodi}{00xx0xxxxxxxx}

% ── Dekan Fakultas ───────────────────────────────────────────
% Default di bawah adalah Dekan Fakultas Ilmu Komputer.
% Sesuaikan bila Fakultas Bisnis dan Ilmu Sosial.
\newcommand{\NamaDekan}{Prof. Dr. Eng. Ir. Imam Tahyudin, M.M.}
\newcommand{\NIKDekan}{2012.09.1.009}

% ── Program Studi & Fakultas ─────────────────────────────────
% Fakultas: "Fakultas Ilmu Komputer" (Prodi SI, TI, Informatika)
%           atau "Fakultas Bisnis dan Ilmu Sosial" (Bisnis Digital,
%           Ilmu Komunikasi). Uppercase diterapkan otomatis di cover.
\newcommand{\NamaProdi}{Informatika}
\newcommand{\NamaFakultas}{Fakultas Ilmu Komputer}

% ── Gelar yang Diperoleh ─────────────────────────────────────
\newcommand{\GelarSkripsi}{Sarjana Komputer (S.Kom.)}

% ── Tahun & Tanggal ──────────────────────────────────────────
\newcommand{\TahunUjian}{2026}            % tahun ujian skripsi
\newcommand{\TanggalACC}{1 Januari 2026}  % tanggal ACC terakhir dosen pembimbing
\newcommand{\TanggalUjian}{20 Januari 2026} % tanggal ujian/sidang skripsi
\newcommand{\TanggalKataPengantar}{Purwokerto, \today}  % penulisan kata pengantar
\newcommand{\KotaPernyataan}{Purwokerto, \TanggalUjian} % tempat/tanggal surat pernyataan

% ── Kata Kunci Abstrak (3–5 kata) ────────────────────────────
% Default di bawah = placeholder template asli (Template_Skripsi.docx).
\newcommand{\KataKunci}{satu, dua, tiga, empat, lima}
\newcommand{\Keywords}{satu, dua, tiga, empat, lima}
